\documentclass[12pt,a4paper]{article}
\usepackage[top=2.5cm,bottom=2.5cm,left=2.5cm,right=2.5cm,headheight=15pt]{geometry}
\usepackage{amsmath,amssymb}
\usepackage{tikz}
\usetikzlibrary{shapes,arrows,arrows.meta,positioning,calc,backgrounds,fit,matrix}
\usepackage{booktabs,array,longtable,tabularx,multirow}
\usepackage{xcolor}
\usepackage{enumitem}
\usepackage{fancyhdr}
\usepackage{titlesec}
\usepackage{mdframed}
\usepackage{tcolorbox}
\tcbuselibrary{skins,breakable}
\usepackage{hyperref}

%% ---- Colors ----
\definecolor{folblue}{RGB}{0,82,165}
\definecolor{lightblue}{RGB}{220,235,255}
\definecolor{darkorange}{RGB}{180,70,0}
\definecolor{lightorange}{RGB}{255,235,210}
\definecolor{darkgreen}{RGB}{20,110,20}
\definecolor{lightgreen}{RGB}{200,240,205}
\definecolor{darkred}{RGB}{150,20,20}
\definecolor{lightred}{RGB}{255,210,210}
\definecolor{lightgray}{RGB}{242,242,242}

%% ---- tcolorbox styles ----
\newtcolorbox{qbox}[2][]{enhanced,breakable,
  colback=lightorange,colframe=darkorange,arc=4pt,
  title={\textbf{#2}},fonttitle=\bfseries\color{white},
  attach boxed title to top left={yshift=-2mm,xshift=4mm},
  boxed title style={colback=darkorange,arc=3pt},#1}

\newtcolorbox{solbox}[2][]{enhanced,breakable,
  colback=lightblue,colframe=folblue,arc=4pt,
  title={\textbf{#2}},fonttitle=\bfseries\color{white},
  attach boxed title to top left={yshift=-2mm,xshift=4mm},
  boxed title style={colback=folblue,arc=3pt},#1}

\newtcolorbox{keybox}[2][]{enhanced,breakable,
  colback=lightgreen,colframe=darkgreen,arc=4pt,
  title={\textbf{#2}},fonttitle=\bfseries\color{white},
  attach boxed title to top left={yshift=-2mm,xshift=4mm},
  boxed title style={colback=darkgreen,arc=3pt},#1}

\newtcolorbox{algobox}[2][]{enhanced,breakable,
  colback=lightgray,colframe=gray!60,arc=4pt,
  title={\textbf{#2}},fonttitle=\bfseries,
  attach boxed title to top left={yshift=-2mm,xshift=4mm},
  boxed title style={colback=gray!60,arc=3pt},#1}

\newtcolorbox{warnbox}[2][]{enhanced,breakable,
  colback=lightred,colframe=darkred,arc=4pt,
  title={\textbf{#2}},fonttitle=\bfseries\color{white},
  attach boxed title to top left={yshift=-2mm,xshift=4mm},
  boxed title style={colback=darkred,arc=3pt},#1}

%% ---- Page style ----
\pagestyle{fancy}\fancyhf{}
\lhead{\color{folblue}\textbf{Hash Functions \& Hash Table Collision Resolution --- CSE 717 InfoSec}}
\rhead{\color{folblue}Exam: 17 Jun 2026}
\cfoot{--- \thepage\ ---}
\renewcommand{\headrulewidth}{0.4pt}

%% ---- Section style ----
\titleformat{\section}{\large\bfseries\color{folblue}}{}{0em}{}[\titlerule]
\titleformat{\subsection}{\normalsize\bfseries\color{darkorange}}{}{0em}{}

\begin{document}

\begin{center}
{\LARGE\bfseries\color{folblue} CSE 717 --- Information Security}\\[4pt]
{\large\bfseries Topic 4: Hash Functions \& Hash Table Collision Resolution}\\[2pt]
{\large Complete Past Paper Solutions: 2020--2024}\\[6pt]
{\normalsize Source: \textit{LcMaterial\#2\_InfoSecuritySP26.pdf} (Hashing, pp.2--25)}\\[2pt]
{\normalsize\color{darkred}\textbf{Every Topic-4 question 2020--2024 included. Zero omissions.}}
\end{center}

\tableofcontents
\newpage

%% ==========================================
\section*{Quick Reference --- Hashing Cheat Sheet}
%% ==========================================

\begin{keybox}{What hashing is for (Lc\#2 pp.2--10)}
\begin{itemize}[noitemsep]
  \item \textbf{Hash table:} stores key$\to$value pairs (like labelled lockers --- need the right key to access a slot).
  \item \textbf{Hashing is irreversible} (unlike encryption, which is reversible) --- used for \textbf{security, integrity, authentication}, NOT for hiding data.
  \item \textbf{Applications:} (1) Password storage --- store hash(password), compare hashes on login; (2) Data integrity --- compare file hash before/after transfer; (3) Digital signatures --- sign the hash with private key; (4) MAC/HMAC --- verify message integrity + sender authenticity; (5) Blockchain --- block linking, Proof-of-Work, tamper resistance.
\end{itemize}
\end{keybox}

\begin{keybox}{Properties of a Good Hash Function (6, Lc\#2 pp.11--13)}
\begin{enumerate}[noitemsep]
  \item \textbf{Deterministic:} same input $\to$ always the same output hash.
  \item \textbf{Quick Computation:} hash must be computable fast, or the system isn't efficient.
  \item \textbf{Pre-Image Resistance:} given $H(A)$, infeasible to recover $A$.
  \item \textbf{Small change $\to$ big hash change (avalanche):} a tiny change in input causes a huge change in the output hash.
  \item \textbf{Collision Resistant:} infeasible to find $A \neq B$ with $H(A)=H(B)$.
  \item \textbf{Puzzle Friendly:} hard to pick an input that produces a pre-defined output --- inputs should be selectable from as wide a distribution as possible.
\end{enumerate}
\end{keybox}

\begin{algobox}{Collision-Resolution Methods (Open Addressing)}
\begin{itemize}[noitemsep]
  \item \textbf{Linear Probing:} if slot $h(k) \bmod m$ is full, try $(h(k)+1)\bmod m$, then $+2$, $+3$, \dots --- check the \emph{next} slot sequentially.
  \item \textbf{Quadratic Probing:} if $h(k)\bmod m$ is full, try $(h(k)+1^2)\bmod m$, then $(h(k)+2^2)\bmod m$, then $(h(k)+3^2)\bmod m$, \dots
  \item \textbf{Double Hashing:} use a 2nd hash function $h_2(k)$ for the step size --- probe sequence is
  \[
  h_{\text{new}} = \big(h_1(k) + i \cdot h_2(k)\big) \bmod m, \quad i = 0, 1, 2, \dots
  \]
  (most collision-resistant of the three --- step size varies per key, avoids clustering).
\end{itemize}
\end{algobox}

\newpage
%% ==========================================
\section{2024 --- Section B Q7(c): Double Hashing Insertion (3 marks)}
%% ==========================================

\begin{qbox}{Question --- 2024 Section B Q7(c)}
Insert the keys $8, 47, 22, 44, 39, 32$ (in this order) into a hash table of size $m=13$ using \textbf{double hashing}, with $h_1(k) = k \bmod 13$ and $h_2(k) = 7 + (k \bmod 7)$. Show the resultant hash table. \hfill [3]
\end{qbox}

\begin{solbox}{Answer --- 2024 Section B Q7(c)}
\textbf{Step 1 --- Compute $h_1(k)$ and $h_2(k)$ for each key:}
\begin{center}
\renewcommand{\arraystretch}{1.2}
\begin{tabular}{c|c|c}
\toprule
$k$ & $h_1(k)=k\bmod 13$ & $h_2(k)=7+(k\bmod 7)$ \\
\midrule
8  & 8 & $7+1=8$ \\
47 & 8 & $7+5=12$ \\
22 & 9 & $7+1=8$ \\
44 & 5 & $7+2=9$ \\
39 & 0 & $7+4=11$ \\
32 & 6 & $7+4=11$ \\
\bottomrule
\end{tabular}
\end{center}

\textbf{Step 2 --- Insert in order} (probe sequence $h_{\text{new}}=(h_1(k)+i\cdot h_2(k))\bmod 13$, $i=0,1,2,\dots$):
\begin{itemize}[noitemsep]
  \item \textbf{8:} $h_1=8$, slot 8 empty $\to$ \textbf{place at 8}.
  \item \textbf{47:} $h_1=8$, slot 8 \emph{occupied} (by 8). Try $i=1$: $(8+12)\bmod 13 = 20\bmod13 = 7$, slot 7 empty $\to$ \textbf{place at 7}.
  \item \textbf{22:} $h_1=9$, slot 9 empty $\to$ \textbf{place at 9}.
  \item \textbf{44:} $h_1=5$, slot 5 empty $\to$ \textbf{place at 5}.
  \item \textbf{39:} $h_1=0$, slot 0 empty $\to$ \textbf{place at 0}.
  \item \textbf{32:} $h_1=6$, slot 6 empty $\to$ \textbf{place at 6}.
\end{itemize}

\textbf{Final hash table ($m=13$):}
\begin{center}
\renewcommand{\arraystretch}{1.2}
\begin{tabular}{c|c|c|c|c|c|c|c|c|c|c|c|c|c}
\toprule
Slot & 0 & 1 & 2 & 3 & 4 & 5 & 6 & 7 & 8 & 9 & 10 & 11 & 12 \\
\midrule
Key  & 39 & -- & -- & -- & -- & 44 & 32 & 47 & 8 & 22 & -- & -- & -- \\
\bottomrule
\end{tabular}
\end{center}
\textit{Only key 47 collided (with 8 at slot 8) and needed 1 extra probe via $h_2$.}
\end{solbox}

\newpage
%% ==========================================
\section{2023 --- Section B Q8(c): Linear Probing Insertion (3 marks) + Q5(b): Hash Function Properties (2.5 marks)}
%% ==========================================

\begin{qbox}{Question --- 2023 Section B Q8(c)}
Insert the keys $21, 8, 13, 44, 28, 33$ (in this order) into a hash table of size $m=13$ using \textbf{open addressing with linear probing}, with primary hash function $h(k) = (k+2) \bmod 13$. Show the resultant hash table. \hfill [3]
\end{qbox}

\begin{solbox}{Answer --- 2023 Section B Q8(c)}
\textbf{Step 1 --- Compute $h(k) = (k+2)\bmod 13$ for each key:}
\begin{center}
\renewcommand{\arraystretch}{1.2}
\begin{tabular}{c|c|c}
\toprule
$k$ & $k+2$ & $h(k)=(k+2)\bmod 13$ \\
\midrule
21 & 23 & 10 \\
8  & 10 & 10 \\
13 & 15 & 2 \\
44 & 46 & 7 \\
28 & 30 & 4 \\
33 & 35 & 9 \\
\bottomrule
\end{tabular}
\end{center}

\textbf{Step 2 --- Insert in order} (linear probe: if slot full, try next slot $+1$, wrap around mod 13):
\begin{itemize}[noitemsep]
  \item \textbf{21:} $h=10$, slot 10 empty $\to$ \textbf{place at 10}.
  \item \textbf{8:} $h=10$, slot 10 \emph{occupied} (by 21). Probe slot 11 --- empty $\to$ \textbf{place at 11}.
  \item \textbf{13:} $h=2$, slot 2 empty $\to$ \textbf{place at 2}.
  \item \textbf{44:} $h=7$, slot 7 empty $\to$ \textbf{place at 7}.
  \item \textbf{28:} $h=4$, slot 4 empty $\to$ \textbf{place at 4}.
  \item \textbf{33:} $h=9$, slot 9 empty $\to$ \textbf{place at 9}.
\end{itemize}

\textbf{Final hash table ($m=13$):}
\begin{center}
\renewcommand{\arraystretch}{1.2}
\begin{tabular}{c|c|c|c|c|c|c|c|c|c|c|c|c|c}
\toprule
Slot & 0 & 1 & 2 & 3 & 4 & 5 & 6 & 7 & 8 & 9 & 10 & 11 & 12 \\
\midrule
Key  & -- & -- & 13 & -- & 28 & -- & -- & 44 & -- & 33 & 21 & 8 & -- \\
\bottomrule
\end{tabular}
\end{center}
\textit{Only key 8 collided (with 21 at slot 10) and needed 1 extra linear probe to slot 11.}
\end{solbox}

\subsection{Part Q5(b): Properties of a Good Hash Function}
\begin{qbox}{Question --- 2023 Section B Q5(b)}
Enumerate the properties of a good hash function. \hfill [2.5]
\end{qbox}
\begin{solbox}{Answer --- 2023 Section B Q5(b)}
\begin{enumerate}[noitemsep]
  \item \textbf{Deterministic:} same input always produces the same hash output.
  \item \textbf{Quick Computation:} the hash must be computed efficiently/quickly for any input.
  \item \textbf{Pre-Image Resistance:} given $H(A)$, it is infeasible to find $A$.
  \item \textbf{Avalanche effect:} a small change in input causes a large, unpredictable change in the output hash.
  \item \textbf{Collision Resistant:} infeasible to find two different inputs $A \neq B$ with $H(A)=H(B)$.
  \item \textbf{Puzzle Friendly:} infeasible to choose an input that yields a pre-defined output --- outputs appear uniformly distributed.
\end{enumerate}
\textit{For 2.5 marks, naming + one-line definition of 5 of these 6 is sufficient --- pick the first 5 if pressed for time.}
\end{solbox}

\newpage
%% ==========================================
\section*{Exam Strategy Summary}
%% ==========================================

\begin{keybox}{High-Yield Checklist for Hash Functions \& Collision Resolution (2/5 years, 2023+2024, 3 marks each)}
\begin{enumerate}[noitemsep]
  \item \textbf{Both 2023 and 2024 ask a hash-table insertion trace} --- 2023 used \textbf{linear probing} (open addressing, simple $h(k)=(k+c)\bmod m$), 2024 used \textbf{double hashing} ($h_1+i\cdot h_2 \bmod m$). \textbf{Practice BOTH} on a fresh 6-key list with $m=10$--$13$ before the exam.
  \item \textbf{Mechanical drill:} for each key (in order) --- compute primary hash, check if slot is free; if occupied, apply the probe sequence (linear: $+1,+2,\dots$; double hashing: $+i\cdot h_2(k)$) until an empty slot is found.
  \item \textbf{Hash function properties} (2023 Q5b, 2.5 marks): memorize the 6-term list --- Deterministic, Quick Computation, Pre-Image Resistance, Avalanche, Collision Resistant, Puzzle Friendly. Cheap, pure-recall marks.
  \item \textbf{Quadratic probing} ($f(i)=i^2$) appears in Lc\#2 with worked examples but NOT in 2020--2024 papers --- know the formula as a backup, don't over-invest.
  \item \textbf{Don't confuse the three methods:} linear probing = fixed step +1; quadratic probing = step $i^2$; double hashing = step $i\cdot h_2(k)$ (a second, key-dependent hash function).
\end{enumerate}
\end{keybox}

\end{document}
